\subsection{Discrete series and the local theta correspondence}\label{ss:B-discrete-series}
At an infinite place the unitary group of a hermitian space of signature $(p,q)$ is $\U(p,q)$. Its discrete series
are indexed by Harish-Chandra parameters $\lambda$, regular vectors in $\frac{n-1}2+\Z^n$ (here written doubled,
$l=2\lambda$, on the coordinates of a space with signs $s$), and the \emph{lowest $K$-type} of $\lambda$ has highest
weight $\Lambda=\lambda+\rho-2\rho_c$. The forms of \emph{type} $l$ are those on which the real group acts through
the minimal idempotent of the discrete series with parameter $l$ (the \emph{kernel}); a product of such kernels over
the infinite places is the \emph{type kernel}. If $l$ has gap $h$, the kernel is integrable with
matrix coefficients decaying like $\exp(-h\cdot\mathrm{dist})$, and integration against it kills bounded functions
invariant under the unipotent radical of a parabolic (Godement's argument \cite[Ch.~7]{Wallach1988}); this is the
statement that the type kernel is a \emph{cusp kernel} of order $h/2$.
\lean{New.HCParam.lowestKType, New.IsHCParam, New.hcParameter, New.highestVector, New.typeKernel, New.IsCuspKernel, New.DiscreteSeries.integral_idempotentKernel_mul_eq_zero, New.DiscreteSeries.exists_conj_mem_dilationGroup}

For a dual pair $(\U(W),\U(V))$ in the Fock model, a polynomial in the variables $\Hom(W,V)$ is a \emph{joint
harmonic} for parameters $(l,l')$ if it is nonzero, harmonic along rows and columns, and highest for both groups of
the weights of the two lowest $K$-types. A datum is an \emph{arrow} from type $l$ to type $l'$ if its lifts take forms of
type $l$ to forms of type $l'$, and an \emph{edge} if it is an arrow in both directions.

\begin{theorem}[Theta correspondence for discrete series]\label{thm:B-theta-ds}
If $l$ and $l'$ satisfy the block rule (\cref{ss:centres}), then they \emph{match}: every row and column is matched
with a sign condition or lies in the central range left by the difference of the ranks. For matching parameters a
joint harmonic exists, given as a product of minors along layers; its projection by the discrete series kernel
$e_\lambda$ has $\re\langle\Phi,\pi(e_\lambda)\Phi\rangle>0$, is fixed by every function reproducing $e_{\lambda'}$, and
generates all projected polynomial vectors under bump functions of the second group. Consequently the Fock vector of a
joint harmonic is an edge between the types $l$ and $l'$.
\end{theorem}
\lean{New.isBlockRule_iff_isMatch, New.IsMatch, New.IsMatch.exists_isJointHarmonic, New.IsJointHarmonic, New.IsJointHarmonic.howe_argument, New.OscillatorPair.re_inner_integratedRep_pos, New.OscillatorPair.integratedRep_integratedRep_eq_self, New.OscillatorPair.exists_bump_integratedRep_eq, New.IsEdge, New.IsArrow, New.FockDualPair.std, New.staircaseWeight}

\begin{proof}[Sources]
The shape of the correspondence of discrete series parameters is due to J.-S.~Li \cite{Li1990} and A.~Paul
\cite{Paul1999}; the joint harmonics are those of Kashiwara and Vergne \cite{KashiwaraVergne1978}, and the generation
statement is Howe's argument \cite{Howe1989}. The Lean development proves the statements used: matching from the block
rule as a statement about integer vectors; the existence of joint harmonics by layers of minors; positivity of the
projection by comparing the Casimir operators of the two groups on a joint harmonic; and Howe's argument in the Fock
model.
\end{proof}

\subsection{Characters of finite general linear groups}\label{ss:B-green}
Let $\mathbb k$ be a finite field and $T$ an injective character of $\bar{\mathbb k}^\times$. For a label $c$ of degree $r$ the
character $\Xi_c$ of $\GL_r(\mathbb k)$ is the parabolic induction product of Green's functions $\tau_P$ over the irreducible
factors $P$ of $c$ \cite{Green1955}.

\begin{theorem}\label{thm:B-green}
The $\Xi_c$, $c$ a label of degree $r$, are orthonormal irreducible characters of $\GL_r(\mathbb k)$; the elements $e_c$ formed
from them are orthogonal central idempotents, and their tensor products over finite families of fields are orthogonal central
idempotents. Multiplying $\Xi_c$ by $\chi\circ\det$ multiplies the roots of $c$ by the element $x_\chi\in\mathbb k^\times$ naming
$\chi$. If $l\ne0$ and $x_\chi$ is not a root of $c$, then $e_c$ kills $\sum\chi(\det h)^{-1}h$ over the invertible $h=1+ul^{\mathsf T}$.
\end{theorem}
\lean{New.greenTau, New.labelChar, New.labelIdem, New.IsLabel.isIrreducibleChar_labelChar, New.labelIdem_orthogonal, New.isCentralProj_tupleIdem, New.tupleIdem_orthogonal, New.evalGL_labelChar_mul_det, New.hyperplaneSum_mul_labelIdem}

\begin{proof}[Sources]
Green's theory of the characters of $\GL_r(\mathbb k)$ \cite{Green1955}: the $\Xi_c$ for squarefree $c$ are the characters
induced from the Coxeter tori of the factors; orthonormality is computed in Hall's algebra. The hyperplane sum spans
representations induced from the stabiliser of $\ker l$ with $\chi\circ\det$ on the line, whose labels have the root $x_\chi$.
\end{proof}

\subsection{Algebraic modular forms and Bessel periods}\label{ss:B-bessel}
For a definite hermitian space $V$ over the CM field $E$, automorphic forms on $\U(V)$ of a given weight and level are
\emph{algebraic modular forms} in the sense of Gross \cite{Gross1999}: functions on the finite set of double cosets
$\U(V)(E^+)\backslash\U(V)(\Ad_{\mathrm f})/K$ with values in the Weyl module of the weight, transforming by the stabilisers. The
Petersson product is the finite sum over the double cosets weighted by the inverse orders of the stabilisers; for free levels
(\cref{prop:free-levels}) all weights are $1$.
\lean{New.algFormRep, New.levelSpace}

For a hermitian space $V$ of rank $n$ with a point, the \emph{first minimum} $\lambda(g)$ at $g$ is the least length of a nonzero vector of the lattice of the moved space at $g$, cut off at $1$, and the
\emph{height} is $1/\lambda(g)$. A \emph{form of type $e$}, for a cusp kernel $e$ (a minimal idempotent for convolution on the real group that decays polynomially and is reproduced by a compactly supported
function), is a bounded automorphic function fixed by right convolution with $e$.

\begin{theorem}[Cusp estimate and rapid decay]\label{thm:B-cusp}
Let $N\ne0$ be a natural number in the different of $E$, $\kappa=\kappa(E,n,N)\ge1$ the cusp constant, $M\ne0$ and $0<\delta\le1$.
\begin{enumerate}[label=\textup{(\alph*)}]
\item If $\lambda(g)<(\delta/(M\kappa^2))^n$, the short lattice vectors at $g$ span a rational isotropic subspace $A\ne0$, and the unipotent radical of its stabiliser is covered by its rational points times the real ball
of radius $\delta$ about $1$ times the principal level $M$.
\item Forms of type $e$ are cuspidal: their constant terms along the unipotent radicals of nonzero rational isotropic subspaces vanish.
\item Let $F$ be a form of type $e$ and principal level $M$. If $\lambda(g)<(\delta/(M\kappa^2))^n$ and $\norm{F(gx)-F(g)}\le C$ for every real $x$ within $\delta$ of $1$, then $\norm{F(g)}\le C$.
\item Forms of type $e$ decay faster than every power of the height, uniformly at a fixed level, and the forms of a fixed type and level span a space of finite dimension.
\end{enumerate}
\end{theorem}
\lean{New.Space.CuspFrame.length_mul_le, New.Space.adelicCoords, New.Space.mem_Gamma_of_adeleMatrix_eq_radicalElem, New.Space.CuspFrame.spans, New.Space.IsCuspidal, New.Space.adelicCoords_gram, New.cuspConst, New.Space.CuspFrame.ne_zero, New.Space.CuspFrame.perpCols, New.Space.CuspFrame.col_isotropicCols_mem, New.Space.IsCuspidal.norm_le_of_firstMinimum_lt, New.Space.exists_cuspFrame_of_height, New.Space.prod_opNorm_mul_le, New.Space.exists_norm_mul_height_pow_le, New.Space.ratEmbed_move, New.Space.exists_cuspFrame, New.exists_isIntBounded_sub_map, New.Space.CuspFrame.isotropic, New.Space.CuspFrame.lengthX_le, New.Space.adelicCoords_evalInf_ne_zero, New.one_le_cuspConst, New.Space.IsForm.mem_rapidFunctions, New.Space.exists_mem_Gamma_adeleMatrix_mul_eq_radicalElem, New.Space.CuspFrame.radicalCovered, New.Space.CuspFrame.orthogonal, New.Space.mem_dilationGroup_of_ratEmbed_mem_radical, New.Space.isCocompact_radical, New.Space.CuspFrame.lengthX, New.Space.exists_norm_sub_le_half, New.Space.CuspFrame.lengthY, New.Space.isIntBounded_adelicCoords, New.Space.norm_le_of_oscillation, New.rank_of_apply_min, New.dist_le_and_mem_principalLevel_of_adeleMatrix_eq, New.Space.finite_levelSpace, New.Space.adelicCoords_injective, New.Space.one_le_norm_mul_of_span_eq_top, New.Space.adelicCoords_gram_eq_zero, New.Space.map_radicalElem, New.exists_orthogonal_one_sub, New.Space.CuspFrame.isotropicCols, New.Reduction.IsIntBounded.mono, New.Space.dist_conj_one_le, New.Space.CuspFrame.col_perpCols_mem, New.IsCuspKernel.integral_mul_eq_zero, New.Space.integral_unipotentRadical_eq_zero, New.Space.IsForm.isCuspidal, New.Space.IsCuspidal.integral_eq_zero, New.Space.nonempty_cuspFrame, New.Space.isForm_realConv, New.Space.fixesLeft_move, New.Space.IsCuspidal.continuous, New.Space.IsForm.norm_le_of_firstMinimum_lt, New.exists_le_rev_and_mul_rev_lt, New.Space.CuspFrame, New.Space.RadicalCovered, New.Space.radical, New.Space.exists_norm_le_of_radicalCovered, New.Space.adelicCoords_mul_add_mul_eq_one, New.Space.exists_basis_prod_norm_le, New.Space.IsCuspidal.rat_mul}

\begin{proof}[Sources]
These are the classical results on cusp forms \cite[\S\S7--8]{Borel2007}, \cite{MillerSchmid2012}, in a form adapted to bounded functions reproduced by an integrable idempotent, with no differential equation and
with an explicit depth. (a) A Minkowski-reduced basis is sorted by length; if the first minimum is below the threshold, the shortest vectors are isotropic and span an isotropic subspace, and the rational points of the
radical are cocompact in it with an explicit fundamental set. (b) Let $F$ be bounded, continuous, fixed by right convolution with $e$ and invariant under the rational points of a unipotent radical $N$ of finite covolume, and let $\Psi(x)$ be the mean of $F(gvx)$
over $v$ in the quotient of $N$. Then $\Psi$ is fixed by convolution with $e$, since $F$ is (exchange the two integrals), and it is killed by that convolution, being left invariant under the real points of
the radical: a cusp kernel annihilates every bounded continuous function left invariant under the real points of the unipotent radical of an isotropic subspace, a case of Harish-Chandra's
theorem that the $K$-finite coefficients of a discrete series have vanishing constant terms. Hence the mean vanishes. (c) By (a), $g$ times the radical lies in rational points times $\delta$-small real points times the level, so all
values of $F$ on $gN$ lie within $C$ of $F(g)$, while their mean is $0$ by (b); hence $\norm{F(g)}\le C$. (d) High in the cusp, (c) with $\delta^n=\kappa/\mathrm{ht}(g)$ and the bounded
variation of a form under small real translations give the decay; elsewhere the height is at most $\kappa$. For the finite dimension, a form of level $K$ exceeds half its bound only at height at most a fixed
$2R$, so only inside $\Gamma C$ with $C$ compact (Mahler's criterion); cover $C$ by finitely many neighbourhoods of points $a^{-1}$ on which such a form stays within half its bound of its value at $a^{-1}$. A form
vanishing at all these points has every bound halved, hence is $0$, so evaluation at finitely many points is injective on the forms of level $K$.
\end{proof}

For $Y_0=\langle e_1,\dots,e_\varrho\rangle\oplus X_0\oplus\langle z\rangle\oplus\langle e^*_1,\dots,e^*_\varrho\rangle$, the \emph{Bessel
subgroup} $R_0=U\rtimes\U(X_0)$ is the stabiliser of the isotropic flag in the $e_a$ that fixes $z$ modulo the flag; its unipotent
radical $U$ carries the generic character $\psi_U(u)=\psi(\sum_a\langle ue_{a+1},e^*_a\rangle+\langle uz,e^*_\varrho\rangle)$. Every
element of the Bessel group is a product $uh$ with $u\in U$ and $h\in\U(X_0)$. These are the Bessel models of Gan, Gross and Prasad
\cite{GGP2012}; Bessel functions on the real group are the bounded functions transforming by $\psi_U$ on the left and by a
representation of $\U(X_0)$.
\lean{New.Grading, New.Shape.Form.besselS, New.Shape.Form.besselU, New.exists_besselU_mul_toBesselH, New.IsBesselFunction}

\subsection{Constant terms and the unfolding of Bessel periods}\label{ss:B-unfolding}
Let $V$, $V'$ be graded hermitian spaces with strings of lengths $s$ and $s-1$ (\cref{ss:graded}), $\Phi$ a datum on $\Hom(V,V')$ and
$\omega_\alpha$ the Weil representation of $\U(V)\times\U(V')$ twisted by characters $\alpha$ of $\U(1)(\Ad)$. The theta kernel
$\theta(\omega_\alpha(r,h)\Phi)$ is expanded along the rational maps $V\to V'$ modulo those of degree $\ge0$: each class contributes a
constant term, the mean of the theta function along the maps of positive degree, which is a theta value of a zero term in a
mixed model of the Weil representation \cite{Weil1964}. Over a set of classes stable under the rational points of the Bessel group
in which no chosen representative agrees with the shift on the entries of positive grade, the terms integrate to zero against a
$\psi_U$-equivariant function; the classes that survive are the translates of the shift by the rational points of $\U(V')$.

\begin{proposition}[One unfolding step]\label{prop:B-unfolding}
Let $W$ be a continuous $\psi_U$-equivariant function on the Bessel group $R_V$, left invariant under $R_V(E^+)$, and $y$ a rapidly
decreasing automorphic function on $\U(V')(\Ad)$. Put $\mathcal W'(h)=\int_{[R_V]}W(r)\,C_{\mathrm{sh}}(\Phi)(r,h)\,dr$. Then
$\int_{[R_V]}W(r)\,\theta(\omega_\alpha(r,h)\Phi)\,dr=\sum_\gamma\mathcal W'(\gamma h)$, the sum over $R_{V'}(E^+)\backslash\U(V')(E^+)$, and
\[
  \int_{[R_V]}W(r)\int_{[\U(V')]}\theta(\omega_\alpha(r,h)\Phi)\,y(h)\,dh\,dr=\int_{R_{V'}(E^+)\backslash\U(V')(\Ad)}\mathcal W'(h)\,y(h)\,dh,
\]
both sides converging absolutely. If the constants of the two gradings agree, $\gamma\mapsto\mathcal W'(\gamma h)$ is again
$\psi_U$-equivariant, and $\mathcal W'$ is continuous and left invariant under $R_{V'}(E^+)$.
\end{proposition}
\lean{New.Grading.hasSum_unfoldStep, New.Grading.integral_besselGroup_unfold, New.Grading.continuous_unfoldStep, New.Grading.unfoldStep_rat_mul, New.Grading.isEquivariant_unfoldStep, New.Grading.hasSum_constTerm, New.Grading.constTerm_mul, New.Grading.constTerm_eq_theta_zeroTerm}

\begin{proof}[Sources]
This is the transfer of Bessel and Fourier--Jacobi periods under one theta lift, in the unitary case with a definite complement;
see Gan and Ichino \cite{GanIchino2016} for the local theory and Kudla \cite{Kudla1984} for the seesaw principle.
\end{proof}

Two further facts on graded lattices are used. \emph{Iwasawa decomposition for the flag of the grades}: if $E$ contains a primitive root of unity
of order $5040$, every finite point $x$ has a unipotent $u$ such that the lattice of $ux$ is graded. \emph{Finiteness}: for an
invertible idele $\mathfrak q$ of valuation $1$ outside a finite set $Q$, the double cosets under the finite Bessel group and the
lattice stabiliser whose orbits contain a point with the lattice bounds have finitely many representatives, supported on
$Q\cup\bar Q$; this is the finiteness of the lattices between two given ones, a lattice being determined by its local components
\cite[\S81]{OMeara1973}.
\lean{New.Grading.exists_isGraded, New.Grading.exists_finite_representatives, New.Grading.finite_image_latticeMinus, New.Grading.HasLatticeBounds.latticeMinus_componentOn}

\subsection{Hecke operators at split places and the unramified correspondence}\label{ss:B-hecke-transfer}
Let $w$ be a split place of norm $q$, so that $\U(V)(E^+_w)\cong\GL_n(E_w)$. For $j\le n$, $T_{w,j}$ is the Hecke operator of the double coset of
$\mathrm{diag}(\varpi1_j,1_{n-j})$ under $\GL_n(\OO_w)$, acting on the vectors spherical at $w$, and the \emph{Hecke polynomial} is
$H_w(X)=\sum_{j\le n}(-X)^jq^{\binom j2}T_{w,j}$. On spherical vectors $T_{\bar w,j}$ is the operator of the inverse double coset, the operators at
different split places commute, and in a unitary representation a finite-dimensional stable space of spherical vectors is spanned by joint
eigenvectors.

Let $S$, $T$ be hermitian spaces of ranks $n$ and $n+k$ and let the characters of the Weil representation of $\U(S)\times\U(T)$ be unramified at
$w$. On the vacuum at $w$ the operators $T_{w,j}$ for $T$ and $T_{\bar w,j}$ for $S$ count the subspaces of the kernels of an integral matrix and of
its transpose modulo $w$, and the $q$-binomial theorem gives
\[
  H^{T}_w(X)\,\mathrm{vac}=\prod_{l<k}(1-q^lc\,X)\cdot H^{S}_{\bar w}(cX)\,\mathrm{vac},
\]
with $c=\chi(\varpi_w)q^{n/2}$ for the character $\chi$ of the pair. Through equivariant maps and through the theta pairing this identity passes to
theta lifts: the Hecke polynomial of the lift of a spherical form is that of the form, rescaled, times the $q$-Pochhammer symbol. This is the
unramified theta correspondence for $\GL_n$, stated for operators (Howe; Rallis \cite{Rallis1982}; Kudla \cite{Kudla1986}).
\lean{New.heckeSplit, New.heckeSplit_star, New.heckeSplit_dual, New.heckeSplit_adjoint, New.heckeSplit_comm, New.IsUnitary.le_span_joint_eigenvectors, New.heckePolyEval, New.heckePolyOp, New.HasHeckePolynomial, New.SiegelWeil.heckePolyOp_pairWeilTarget, New.SiegelWeil.hecke_transfer_of_equivariant, New.SiegelWeil.hecke_transfer_of_pairing, New.SiegelWeil.lift_swap_heckePolyOp, New.SiegelWeil.hasHeckePolynomial_formLift, New.SiegelWeil.hasHeckePolynomial_qPochhammer, New.SiegelWeil.qPochhammer}

\begin{theorem}[Chains of Satake roots]\label{thm:B-chains}
Let $\rho$ be a smooth representation of $\GL_N(E_w)$ with an invariant scalar product, $q$ the norm of $w$, and $x\ne0$ a vector fixed by $\GL_N(\OO_w)$ with Hecke polynomial $P=\sum_{j\le N}(-1)^jq^{\binom j2}t_jX^{N-j}$.
\begin{enumerate}[label=\textup{(\alph*)}]
\item If $P(c)=0$ and $\abs c^2q^k=q^N$, then $c,cq,\dots,cq^{k-1}$ are roots of $P$; if $P(c)=0$ and $P(qc)\ne0$, then $\abs c^2\ge q^{N-2}$.
\item Let $U$ be the Hecke operator of $\mathrm{diag}(\varpi,1,\dots,1)$ for the pro-$p$ Iwahori subgroup. If $d\ne0$ and $P(X/d)=\prod_{i\le k}(X-q^i)\cdot S$ with $S$ dividing $R$, then $x$ is orthogonal to
$\bigl((dX-q^k)R(dX)\bigr)(U)y$ for every $y$ fixed by the pro-$p$ Iwahori subgroup.
\end{enumerate}
\end{theorem}
\lean{New.Patching.HasRootChain, New.Patching.adjoint_aeval_map_starRingEnd, New.Patching.aeval_cofactor_eq_zero, New.Patching.aeval_conj_eq_zero, New.Patching.apply_inv_partialHecke_zero_top, New.Patching.chain_lemma, New.Patching.chain_lemma_of_adjoint, New.Patching.compactMean_apply_inv_partialHecke_one_succ, New.Patching.eq_zero_of_partialHecke_zero_top, New.Patching.isRoot_chain_of_norm, New.Patching.le_norm_sq_of_isRoot_of_not_isRoot_mul, New.Patching.pairing_aeval_heckeCoset_eq_zero, New.Patching.partialHecke_hyperplane_count_dual}

\begin{proof}[Sources]
The Satake parameters of unitary unramified representations of $\GL_N$ come in chains; both statements follow from Tadi\'c's classification of the unitary dual \cite{Tadic1986} and the geometric lemma of Bernstein and Zelevinsky. The Lean
development proves them directly, by two counts of hyperplanes: a root $c$ with $qc$ not a root has $\abs c^2\ge q^{N-2}$; with the symmetry $y\mapsto q^{N-1}/\bar y$ of the roots, forced by unitarity, this gives the chain; and the same counts let
the roots of a chain below the top one be cancelled one at a time against the partial Hecke sums of translates of $x$, which gives the orthogonality.
\end{proof}

\subsection{A comparison of simple trace formulas}\label{ss:B-trace-comparison}
Let $V$ be a hermitian space of rank $n$ over the CM field $E$. For a test function $e\otimes\phi$, with $e$ a cusp kernel on the real group (a minimal convolution
idempotent with polynomial decay) and $\phi$ locally constant of compact support on $G_{\rm f}$, the trace of its action on forms of type $e$ is the geometric side
$\sum_R$ over the characteristic polynomials $R$ of rational points of the orbital integrals. A polynomial $R$ that is irreducible with a root off the unit circle contributes
nothing (the Selberg principle for idempotents of the discrete series, after Harish-Chandra); for $R$ irreducible with all roots on the unit circle the orbital integral
splits into a product of characters of the compact unitary group and a factor that depends only on $R$ and $\phi$, which matches the classes of different signs.
\lean{New.Space.orbitalIntegral_eq_zero, New.Space.geometricSide_singleton_eq_zero_of_norm_ne_one, New.Space.orbitalIntegral_eq_prod_mul, New.trace_comparison}

\begin{theorem}[Comparison of traces]\label{thm:B-trace-comparison}
Let $E$ contain the $5040$th roots of unity and have an even number of infinite places, and let $V$, $V'$ be hermitian spaces of rank $n\ge2$ over $E$ with signs $s$, $s'$.
Let $l$ be a doubled Harish-Chandra parameter with gap at least $500n^3+10n^5$ at every infinite place for both signs, $K$ a level, and $a$ an element of the group algebra of
$G_{\rm f}$ preserving the $K$-fixed vectors in every representation, such that $a\star\mathbf 1_K$ is cuspidal at a split place. Then $a$ has the same trace on the forms of
parameter $l$ and level $K$ on $V$ and on $V'$.
\end{theorem}
\lean{New.trace_eq_of_isCuspidalAt, New.isCuspidalAt_convIndicator_tupleIdem, New.sum_coeff_tupleIdem_radicalHom, New.sum_coeff_labelIdem_radicalHom, New.finMean_labelChar_radicalHom}

\begin{proof}[Sources]
The statement for test functions supported, at a second place, on elements with a simple characteristic root $0$ is the simple trace formula with a cuspidal
component, as in Arthur and Clozel \cite[Ch.~1]{ArthurClozel1989} and Harris and Labesse \cite[Prop.~1.10.5]{HarrisLabesse2004}; the second condition is removed by
cutting the level at an auxiliary split place down to the stabiliser of a line, where a line Hecke operator forces the simple zero and counts the spherical vectors through
its determinant. Here integrable idempotent kernels of discrete series with a large gap replace pseudo-coefficients, and there is neither truncation nor endoscopy. The cuspidal
input is Green's theory: the class function of a monic irreducible polynomial of degree $n$ has mean zero on the unipotent radical of every proper coordinate parabolic of
$\GL_n(\F_q)$ \cite{Green1955}, so the idempotent of a tuple of labels with an irreducible member is cuspidal at the place of that member.
\end{proof}

\subsection{The Siegel--Weil formula and Siegel's formula}\label{ss:B-siegel-weil}
Let $S$ be a hermitian space over $E$ split into $p$ hyperbolic planes and $T$ a hermitian space of rank $m$, both unimodular at every finite place, with the Weil representation $\omega$ of the
pair of unitary groups on the test functions $\mathcal S$ of $\mathrm{Hom}(S,T)$ (\cref{ss:B-weil}). A \emph{Lagrangian} of $S$ is a maximal isotropic subspace; every Lagrangian is spanned by the isotropic
frame of a complete splitting $S=lE^p\oplus kE^p$. The \emph{zero term} of $f$ along a Lagrangian $L$ is the mean of $\theta(\varrho(z)f)$ over the maps $z\colon S\to T$ vanishing on $L$, adelic modulo
rational; it is the value at $0$ in the Schr\"odinger model of a splitting along $L$ and depends only on $L$.
\lean{New.SiegelWeil.CompleteSplitting, New.SiegelWeil.splitWeil, New.SiegelWeil.Lagrangian, New.SiegelWeil.homSlice_zero, New.SiegelWeil.forall_colSpan_dot_eq_zero_iff, New.SiegelWeil.exists_mulVec_add_mulVec_of_orthogonal, New.SiegelWeil.mixedValue, New.SiegelWeil.CompleteSplitting.lagrangian, New.SiegelWeil.Lagrangian.eq_of_le, New.SiegelWeil.Lagrangian.exists_completeSplitting, New.SiegelWeil.Lagrangian.completeSplitting, New.SiegelWeil.Lagrangian.numPlanes, New.SiegelWeil.Lagrangian.chosenSplitting, New.SiegelWeil.Lagrangian.lagrangian_chosenSplitting, New.SiegelWeil.Lagrangian.zeroTerm, New.SiegelWeil.zeroTerm_lagrangian_eq_theta, New.SiegelWeil.isometrySplitting, New.SiegelWeil.zeroTerm_isometrySplitting, New.SiegelWeil.zeroTerm_lagrangian_isometrySplitting_comp, New.SiegelWeil.exists_isometry_orthogonal, New.SiegelWeil.mem_comap_iff_forall_dot_eq_zero, New.SiegelWeil.exists_lagrangian_equiv_containing, New.SiegelWeil.mixedValue_splitWeil_levi, New.SiegelWeil.mixedValue_splitWeil, New.SiegelWeil.hasSum_mixedValue_theta_splitWeil, New.SiegelWeil.normalisedTorus, New.SiegelWeil.normalisedTorus_zero, New.SiegelWeil.norm_mixedValue_normalisedTorus, New.SiegelWeil.theta_normalisedTorus_le, New.SiegelWeil.norm_zeroTerm_parabolic_mul, New.Space.Splitting.exists_slice_omega_vacElem_of_isEmpty, New.SiegelWeil.norm_zeroTerm_vacElem_pos, New.SiegelWeil.exists_ideleAbs_le_of_isCompact, New.SiegelWeil.exists_zeroTerm_normalisedTorus_le, New.SiegelWeil.prod_mul_prod_lt_prod_one_add_pow, New.SiegelWeil.transferWeight, New.SiegelWeil.transferPoint, New.SiegelWeil.norm_qPochhammer_mul_norm_eq_norm_sum, New.SiegelWeil.norm_qPochhammer_neg, New.SiegelWeil.norm_transferWeight, New.SiegelWeil.sum_norm_transferWeight_lt, New.SiegelWeil.exists_zeroTerm_omega_schrodingerOp_eq_smul, New.SiegelWeil.exists_zeroTerm_transferPoint_eq, New.SiegelWeil.sum_norm_mul_norm_zeroTerm_le, New.SiegelWeil.eq_zero_of_le_tsum_zeroTerm}

\begin{theorem}[Siegel--Weil formula in Weil's range]\label{thm:B-siegel-weil}
If $m>2p$, the zero terms are summable in operator norm over the Lagrangians of $S$, and for every $f\in\mathcal S$ the function $h\mapsto\theta(\omega(1,h)f)$ is integrable on $[\U(T)]$ with
\[
  \frac1{\vol[\U(T)]}\int_{[\U(T)]}\theta(\omega(1,h)f)\,dh=\sum_L\mathrm{zt}_L(f).
\]
\end{theorem}
\lean{New.SiegelWeil.radicalElement_eq_conj_levi, New.SiegelWeil.exists_compact_parabolic_normal_form, New.SiegelWeil.exists_siegel_set, New.SiegelWeil.tsum_div_two_pow_lt_top, New.SiegelWeil.exists_mul_k_apply_ne_zero, New.SiegelWeil.lagrangian_isometrySplitting_comp_eq, New.SiegelWeil.exists_summable_zeroTerm_bound, New.SiegelWeil.lCoord_mul_levi_dyadicDiag, New.SiegelWeil.tendsto_integral_theta_zeroTerm_sub, New.SiegelWeil.tendsto_zeroTerm_of_notMem, New.SiegelWeil.bddAbove_norm_sub_tsum_norm_zeroTerm, New.SiegelWeil.exists_lagrangian_equiv_col_mem, New.SiegelWeil.tendsto_tsum_zeroTerm_sub_integral_theta, New.SiegelWeil.integrable_theta_splitWeil, New.SiegelWeil.siegel_weil}

\begin{proof}[Sources]
This is Weil's theorem \cite[\S52, Thm.~5]{Weil1965} in the convergent range, with the Eisenstein series written as the sum of the zero terms over the Lagrangians. The proof in the Lean development
follows the classical route: reduction theory gives Siegel sets, every adelic point being $\gamma a_x\kappa$ with $\gamma$ rational, $a_x$ a dyadic torus element and $\kappa$ in a compact set; theta is
bounded uniformly along the normalised torus when $m>\operatorname{rank}S$, and the zero terms of a second splitting are bounded along it. Let $\ell$ be the difference of the two sides. Splitting off
one isotropic line and inducting on the rank of $S$, $\ell$ along the normalised torus tends to $0$ as any coordinate of the torus tends to infinity, so $\ell(\omega(g)f)$ exceeds the sum of the norms of
the zero terms of the translated vacuum by a bounded amount; a continuous invariant functional dominated in this way vanishes when $\operatorname{rank}T>\operatorname{rank}S>0$, by the Hecke maximum principle.
\end{proof}

For a function $\Phi_0$ of finite adelic vectors of rank $n$ and $\alpha\in E^+$ the \emph{density} is $\delta(\Phi_0,\alpha)=\iint\Phi_0(x)\psi_{\rm f}(z(x^*x-\alpha)/2)\,dx\,dz$, over the finite vectors $x$ and the
self-adjoint finite adeles $z$, with the codifferent of volume $1$. For a level $K$ the \emph{class sum} of a function is its sum over the classes, weighted by $1/\#\mathrm{Aut}$; the class sum of $1$ at the
stabiliser of $\hat\OO_E^n$ is the \emph{mass} $\mathrm{mass}_n$.
\lean{New.mem_stdRatPoints_iff, New.classSum_stdRatPoints_eq_mass, New.mass_pos, New.normOneLatticeVectors, New.image_normOneLatticeVectors_eq_orbit_inter, New.classSum_card_norm_one_eq_mass, New.idealZeta, New.le_two_of_pow_le_ofReal, New.lintegral_le_idealZeta_mul, New.coe_stdLattice, New.characterKernel, New.isOpen_characterKernel, New.conductorIdeal, New.cardQuot_conductorIdeal_ne_zero, New.gauss_sum_sq_le, New.lintegral_gauss_sum_le_idealZeta, New.idealZeta_intAdeles_pow_le}

\begin{theorem}[Siegel's formula]\label{thm:B-siegel-formula}
Let $E$ carry Weil data and $n\ge3$.
\begin{enumerate}[label=\textup{(\alph*)}]
\item (One vector.) For the standard definite space of rank $n$, the mean over the unitary quotient of $\sum_\xi\Psi(h^{-1}\xi)$ is $\Psi(0)$ plus the sum over $\beta\in E^+$ of the integrals
$\int\Psi(x)\psi(\beta x^*x/2)\,dx$; for $\alpha\ne0$ the sum over the rational $\xi$ with $\xi^*\xi=\alpha$ has mean $\abs{D_E}^{1/4}\int P$, $P(z)=\int\Psi(x)\psi(z(x^*x-\alpha)/2)\,dx$.
\item (Vectors of one norm.) Let $K$ be a level, $\Phi_0\ge0$ locally constant, $K$-invariant and of compact support, and $\alpha$ totally positive. The class sum of level $K$ of
$\sum_{x^*x=\alpha}\Phi_0(h^{-1}x)$ is the class sum of $1$ times $\abs{D_E}^{1/4}\prod_v\pi^n\alpha_v^{n-1}/(n-1)!$ times $\delta(\Phi_0,\alpha)$.
\item (Gauss sums.) If $\zeta_{\hat\OO}(\mu)\le2$ and $2(\mu+1)\le n$, the integral over $z$ of $\abs{\int_{\hat\OO^n}\psi_{\rm f}(z(w^*w-\alpha)/2)\,dw}$ is at most $\zeta_{\hat\OO}(\mu)\vol(\hat\OO^n)$.
\item (Masses.) If $E$ contains a primitive $5040$th root of unity, the class sum in rank $j+1$ of the number of vectors of norm $1$ is $\mathrm{mass}_j$; and for $j\ge4u+12$ and $8f\le2^u$,
\[
  \Bigl|\log\frac{\mathrm{mass}_{j+1}}{\mathrm{mass}_j}-f\log\frac{\mathrm{rd}_E^{\,j+1/2}\,j!}{(2\pi)^{j+1}}\Bigr|\le\frac2{2^u}.
\]
\end{enumerate}
\end{theorem}
\lean{New.SiegelWeil.shearSplitting, New.SiegelWeil.eq_lagrangian_of_notMem_range_shear, New.SiegelWeil.zeroTerm_shearSplitting, New.SiegelWeil.siegel_weil_one_vector, New.exists_fundamentalDomain_real, New.realAdele, New.exists_isAddFundamentalDomain_realAdele, New.star_adeleMatrix_mulVec_dotProduct, New.integrable_and_integral_eq_tsum_of_fundamentalDomain, New.SiegelWeil.siegel_weil_fourierCoeff, New.normDensity, New.SiegelWeil.exists_mixedValue_eq_gaussian_mul, New.integral_fourierChar_mul_inv_pow, New.integral_fourierChar_mul_gaussian_pow, New.integral_gaussian_mul_stdAddChar, New.classSum_eq_mul_integral, New.classSum_eq_mul_normDensity, New.mass_eq_mass_succ_mul_normDensity, New.norm_sub_one_le_of_normDensity_eq, New.abs_log_mass_ratio_sub_le}

\begin{proof}[Sources]
(a) is the Siegel--Weil formula (\cref{thm:B-siegel-weil}) for the hyperbolic plane against the definite space; the zero terms along the lines of $l+\jmath^{-1}\beta k$, $\beta\in E^+$, are computed in the
Schr\"odinger model, and Laplace's integral for $1/\Gamma$ gives the archimedean factor. (b) Detect the norm $\alpha$ by its Fourier coefficient; this is Siegel's formula \cite{Siegel1935} for hermitian forms
\cite{Braun1941}, read off from Weil's formula for one vector \cite[Thm.~5]{Weil1965}, with the singular series left as one adelic Fourier integral. (c) Sort the characters by their conductor ideal; for
each ideal the Gauss sum is bounded by Weyl differencing, and the sum over the ideals is the zeta value. (d) Apply (b) to the vectors of norm $1$ in the standard lattice and bound the density by (c); this
is an approximate form of the mass formula for unitary groups \cite{GanHankeYu2001}, which suffices here because only ratios of masses in large rank are used.
\end{proof}

The local densities at a finite set of finite places, read as one ring $A$ with integers $R$, are controlled by the same Gauss sums. The classes modulo $gL$ of the lattice $L$ of integral vectors are
compact and open of volume $\abs g^{\#\kappa}$. For a unimodular hermitian form of rank at least $2\mu+1$ and $\zeta_R(\mu)\le2$, a class modulo $gL$ has mass at most $\zeta_R(\mu)\abs g^{-1}$ times its
volume on each quadric, and on an integral quadric the mass of $L$ differs from $1$ by at most $\zeta_R(\mu)-1$; by Fourier inversion the local density of a continuous $f\ge0$ of compact support at
$\delta^{-1}t$ is the integral of $f$ over the quadric, whenever the outer integral converges absolutely, and its modulus is always at most this integral. At a finite place the standard form of rank at
least two has a basis beginning with a hyperbolic pair, so the volume disintegrates along the form into the fibre measures used in \cref{prop:top-floor}(a).
\lean{New.Floors.IntegralPlace.integers, New.Floors.IntegralPlace.data_normalised, New.Floors.IntegralPlace.volume_setOf_mul_mem, New.Floors.IntegralPlace.exists_ideal_cardQuot_eq, New.Floors.IntegralPlace.skewMeasure_setOf_mul_mem, New.Floors.IntegralPlace.isOpen_isCompact_volume_intVectors, New.Floors.IntegralPlace.modulus_le_volume, New.Floors.IntegralPlace.modulus_le_one, New.Floors.IntegralPlace.unitaryLevel_open_compact, New.Floors.IntegralPlace.modulus_det_eq_one, New.Floors.IntegralPlace.exists_content, New.Floors.IntegralPlace.contentIndex, New.Floors.IntegralPlace.measurable_contentIndex, New.Floors.IntegralPlace.lintegral_contentIndex_pow_le, New.Floors.IntegralPlace.topEichlerMatrix, New.Floors.IntegralPlace.gram_topEichlerMatrix, New.Floors.IntegralPlace.topEichler, New.Floors.IntegralPlace.coe_topEichler, New.Floors.IntegralPlace.topEichler_mem_floorGroup, New.Floors.IntegralPlace.orbitPoint_topEichler, New.Floors.IntegralPlace.topEichler_single, New.Floors.IntegralPlace.measurable_topEichler, New.Floors.IntegralPlace.continuousAt_topEichler, New.Floors.IntegralPlace.intUnits, New.Floors.IntegralPlace.measurableSet_unitaryLevel, New.Floors.IntegralPlace.measurableSet_exists_mul_mem_level, New.Floors.IntegralPlace.ascConv_mconv_unfolding, New.Floors.IntegralPlace.psi_skew_eq_one_iff, New.Floors.IntegralPlace.setIntegral_eq_slab, New.Floors.Place.IsDisintegration.setLIntegral_slab, New.Floors.IntegralPlace.tendsto_average_skew, New.Floors.IntegralPlace.integral_skewMeasure_eq_toReal_lintegral, New.Floors.IntegralPlace.enorm_localDensity_le, New.Floors.integral_eq_mul_of_section}
\lean{New.Floors.IntegralPlace.classGaussSum, New.Floors.IntegralPlace.psi_mul_trace, New.Floors.IntegralPlace.conj_psi_mul_psi_add, New.Floors.IntegralPlace.smul_mem_intVectors_of_psi_eq_one, New.Floors.IntegralPlace.mem_latticeMultiple_mul_of_psi_eq_one, New.Floors.IntegralPlace.enorm_classGaussSum_sq_le, New.Floors.IntegralPlace.integrable_classGaussSum, New.Floors.IntegralPlace.integral_classGaussSum, New.Floors.IntegralPlace.measure_residueClass_le_lintegral}
\lean{New.Floors.IntegralPlace.latticeMultiple, New.Floors.IntegralPlace.mem_latticeMultiple, New.Floors.IntegralPlace.coe_latticeMultiple, New.Floors.IntegralPlace.isCompact_latticeMultiple, New.Floors.IntegralPlace.residueClass, New.Floors.IntegralPlace.isCompact_residueClass, New.Floors.IntegralPlace.isOpen_latticeMultiple, New.Floors.IntegralPlace.volume_latticeMultiple, New.Floors.IntegralPlace.isOpen_residueClass, New.Floors.IntegralPlace.volume_residueClass}
\lean{New.Floors.IntegralPlace.skewMeasure_setOf_mul_mem_sq_le, New.Floors.IntegralPlace.exists_unit_sub_mem, New.Floors.IntegralPlace.lintegral_enorm_classGaussSum_le_idealZeta, New.Floors.IntegralPlace.measure_residueClass_le_idealZeta, New.Floors.IntegralPlace.classGaussSum_one_zero_eq_one, New.Floors.IntegralPlace.one_add_enorm_measure_intVectors_sub_one_le}
\lean{New.Floors.Place.isDisintegration_map_quadricMeasure, New.Floors.card_eq_card_pair_sum, New.Floors.IntegralPlace.exists_hyperbolic_basis, New.Floors.IntegralPlace.isDisintegration_of_hyperbolic_basis, New.Floors.IntegralPlace.exists_kernel_isDisintegration_le}

Strong approximation is used in the following weak form. A continuous function on the adelic unitary group of a space with a hyperbolic plane that is left invariant under the rational points,
right invariant under the points supported on a nonempty set $S$ of finite places and their conjugates and under the integral points trivial on $S$, and whose values on the Levi subgroup $\GL_1\times\U(W_0)$
do not depend on the $\GL_1$ component once it is invariant under the unipotent radical, is constant. The proof combines strong approximation for the two opposite unipotent radicals, Witt's theorem on the
rational points of the Levi factor, and weak approximation with the Iwasawa decomposition \cite{Kneser1965}.
\lean{New.Exchange.exists_torus_mul_quasiReflection_fixes, New.Exchange.exists_torus_mul_rational_fixes, New.Exchange.IsAutFunction.apply_eq_apply_one}

\subsection{The constant term along an isotropic line}\label{ss:B-tower}
This subsection supplies the analytic input of the descent of the exchange identity along hyperbolic planes (\cref{prop:removing-planes}). Let $\nu$ be a character of
$[\U(1)]$ of weight $\mu$, with $\mu(w)>0$ at every infinite place $w$, whose tame part is nontrivial exactly on a nonempty finite set $S$ of finite places that contains no conjugate of one of its
places; such a pair $(\nu,S)$ is a \emph{split tame character}. Its \emph{semilocal character} is $\chi(t)=\nu(t/\bar t)$ on $\prod_{w\in S}E_w$, read as $0$ off the units. With $q_w$ the order of the residue
field at $w$ and $L^S(x,\nu)$ the Euler product of the base change $x\mapsto\nu(x/\bar x)$ without the places of $S\cup\bar S$, put, for $e\ge0$,
\begin{gather*}
  L_*(e)=\Bigl(\prod_{w\in S}q_w\Bigr)^{e^2}\prod_{z=1}^{e}\abs{L^S(z,\nu)}^2,\qquad L_*(-e)=L_*(e)^{-1},\\
  L^\sharp(z)=\abs{L^S(z,\nu)}\prod_{w\in S}q_w^{\,z-1/2},
\end{gather*}
the last for real $z\ge1$, where $L^S$ is given by its Euler product.
\lean{New.Exchange.lValueProd, New.Exchange.lValueProdInt, New.Exchange.lSharp}

Let $W'=W\oplus H$ and $V'=V\oplus H$ be hermitian spaces split by adapted splittings $A$ and $X$ along a hyperbolic plane $H$ (frames orthonormal at the infinite places and integral at the
finite places), and let $x_B$ be the adelic vector spanning the isotropic line of $X$. For an adele $a$ the \emph{line term} $\mathrm{lt}_a$ sends a datum on $\Hom(W',V')$ to a datum on $\Hom(W,V)$:
in the model $\beta=\mathrm{true}$ it is translation by the rank one map $a\,x_B\otimes k_A^*$ followed by the zero terms along the line of $A$ and then along that of $X$; in the model
$\beta=\mathrm{false}$ the dual lines are used and the zero terms are taken in the other order. Write $\mathrm{lt}^\infty_a$ for its archimedean part on the Fock spaces. For real $z$, the \emph{cone
integrand} of a Fock vector $v$ for $(W',V')$ at a pair $p=(\kappa,a)$ of a point $\kappa$ of the maximal compact subgroup of $\U(V')$ and a family $a$ of nonzero complex numbers indexed by the
infinite places is
\[
  \mathrm{ci}_z(v)(\kappa,a)=\prod_w\Bigl(\frac{a_w}{\abs{a_w}}\Bigr)^{2\mu(w)}\,\abs a^{z}\;\mathrm{lt}^\infty_a(\kappa v),\qquad \abs a=\prod_w\abs{a_w}^2,
\]
with $\kappa$ acting by the real Weil representation of weight $\lfloor\operatorname{rk}W'/2\rfloor-\mu$ on $\U(V')$. The \emph{cone integral} $\mathrm{cone}_z(v)$ is its integral over the maximal compact
subgroup and the archimedean ideles $(E\otimes\R)^\times$, both with Haar measure, in the model $\mathrm{true}$; the \emph{dual cone integral} $\mathrm{cone}^*_z(v)$ is the same in the model $\mathrm{false}$.
For an integral matrix $y$ at $S$ the \emph{plane datum} of $v$ and $y$ is the datum on $\Hom(W',V')$ with real part $v$, the standard datum of $y$ composed with the projection to $W$ at $S$, and the
vacuum elsewhere; the \emph{constant term} along $A$ at a point $(u,g)$ of the Levi subgroup $\GL_1\times\U(W)$ is the mean over the unipotent radical of the stabiliser of the line of $A$, modulo its
rational points.
\lean{New.Exchange.coneIntegrand, New.Exchange.cone, New.Exchange.dualCone, New.Exchange.lineTerm, New.Exchange.planeDatum, New.Exchange.planeConstantTerm}

\begin{theorem}[The tower property with ramified data]\label{thm:B-tower}
Let $(\nu,S)$ be a split tame character of weight $\mu$, let $\operatorname{rk}W$ be even and $e=\operatorname{rk}V-\operatorname{rk}W/2$.
\begin{enumerate}[label=\textup{(\alph*)}]
\item ($L$-values.) $L^S(x,\nu)\ne0$ for real $x\ge1$, so $L^\sharp(z)>0$ for $z\ge1$; and for every $z\in\N$
\[
  L_*(z+1)=L^\sharp(z+1)^2L_*(z),\qquad L_*(-z)=L^\sharp(z+1)^2L_*(-z-1).
\]
\item (Cone integrals.) For $z>0$ and a Schwartz vector $v$ the cone integrand is integrable for the product Haar measure, and $\mathrm{cone}_z(v)$ and $\mathrm{cone}^*_z(v)$ are Schwartz vectors.
\item (Unfolding.) For a Schwartz vector $v$, an integral $y$ at $S$ and a Levi point $(u,g)$ with $u=1$ or $\operatorname{rk}V=\operatorname{rk}W$, there is a function $\mathcal Z$, analytic on the real
line, such that the constant term along $A$ at $(u,g)$ of the theta integral of $(W',V')$ on the plane datum is $\mathcal Z(e+1)$, and for real $x\ge1$
\[
  \mathcal Z(x)=\varepsilon\,\abs{\alpha_0}^x\prod_{w\in S}q_w^{\,x-1/2}\,L^S(x,\nu)\,\theta_{W,V}\bigl(\mathrm{cone}_x(v),\mathrm{row}_y\bigr)(g),
\]
and $\mathcal Z(1-x)$ is the same expression with $\mathrm{cone}^*_x(v)$ and a constant $\varepsilon'$. Here $\alpha_0$ is the normalising idele of the Weil data, and $\abs\varepsilon$, $\abs{\varepsilon'}$ are
positive numbers independent of $(\nu,S)$.
\item (One plane.) Let $z\in\N$. If $e=z$, there are $C>0$, independent of the character, and $\gamma$ with $\abs\gamma=C\,L^\sharp(z+1)$ such that the constant term in (c) is
$\gamma\,\theta_{W,V}(\mathrm{cone}_{z+1}(v),\mathrm{row}_y)(g)$. If $e=-z-1$, the same holds at $u=1$ with $\mathrm{cone}^*_{z+1}(v)$ in place of $\mathrm{cone}_{z+1}(v)$.
\item (Functional equation.) Let $k\in\N$ with $k\le\mu(w)$ for all $w$. The function $x\mapsto L^S(x,\nu)$ on $[1,\infty)$ extends to an analytic function $\Lambda$ of real $x>-k$ with $\Lambda(j)\ne0$
at every integer $j>-k$, and $L^\sharp(1-z)=C(z)\,L^\sharp(z)$ for real $1\le z<k+1$, where $L^\sharp(1-z)$ is formed with $\Lambda$ and $C(z)>0$ does not depend on the character.
\end{enumerate}
\end{theorem}
\lean{New.Exchange.SplitTameChar.L_ne_zero, New.Exchange.SplitTameChar.weight_ne_zero, New.Exchange.baseChange_poissonVariable, New.Exchange.continuousOn_L, New.Exchange.exists_equiv_principalAdeleUnits, New.Exchange.exists_haar_ideleParts, New.Exchange.exists_lineFinFactor_uniformizer_eq, New.Exchange.finTateIntegrand, New.Exchange.ideleAbs_poissonVariable_false, New.Exchange.ideleParts, New.Exchange.integral_finTateIntegrand, New.Exchange.integral_ideleClass_tsum, New.Exchange.lSharp_pos, New.Exchange.lValueProdInt_natCast, New.Exchange.lValueProdInt_neg, New.Exchange.lValueProdInt_neg_natCast, New.Exchange.lValueProdInt_pos, New.Exchange.lValueProdInt_succ, New.Exchange.lValueProd_pos, New.Exchange.lValueProd_succ, New.Exchange.lineFinFactor, New.Exchange.lineFinFactor_congr, New.Exchange.lineFinFactor_mul_finPart, New.Exchange.lineFinFactor_zero, New.Exchange.measurePreserving_prod_parameters, New.Exchange.measureReal_idelesUnitOff_empty_pos, New.Exchange.poissonVariable, New.Exchange.poissonVariable_false, New.Exchange.poissonVariable_true, New.Exchange.principalAdeleUnit, New.Exchange.realIdele_mul_finIdeleIncl_ideleParts, New.Exchange.valuation_of_lineFinFactor_ne_zero}
\lean{New.Exchange.RadicalCoord, New.Exchange.coe_hermitianPart, New.Exchange.constTermAt, New.Exchange.constTermAt_eq_integral, New.Exchange.constTermAt_smul_k, New.Exchange.constTermAt_zero, New.Exchange.continuous_lineLevi, New.Exchange.eq_one_of_levi_mem_radical, New.Exchange.exists_eq_radicalElement_mul_lineLevi_mul, New.Exchange.exists_integral_lineStabilizer, New.Exchange.exists_radicalElement_eq_mul_lineLevi, New.Exchange.exists_unfold_constant_term, New.Exchange.hasSum_constTermAt, New.Exchange.hermitianPart, New.Exchange.integral_thetaIntegral_lineRadical, New.Exchange.isClosedEmbedding_radicalPoint, New.Exchange.isClosed_radicalSubgroup, New.Exchange.isCocompact_lineStabilizer_subgroupOf, New.Exchange.isHaarMeasure_radicalHaar, New.Exchange.lineLevi_mem_lineStabilizer, New.Exchange.lineStabilizer, New.Exchange.lineSum, New.Exchange.lineSum_smul, New.Exchange.lineSum_splittingLine, New.Exchange.lineTerm_pairWeil_radicalElement, New.Exchange.lintegral_levi_unfold, New.Exchange.lintegral_radicalHaar, New.Exchange.lintegral_radicalHaar_conj_lineLevi, New.Exchange.measurable_lineSum, New.Exchange.mem_lineStabilizer_iff, New.Exchange.mem_radicalSubgroup_iff, New.Exchange.radicalElement_mem_radicalSubgroup, New.Exchange.radicalHaar, New.Exchange.radicalHomeomorph, New.Exchange.radicalMeasure, New.Exchange.radicalPoint, New.Exchange.radicalPoint_hermitianPart, New.Exchange.radicalPoint_mul, New.Exchange.radicalSubgroup, New.Exchange.splittingLine, New.Exchange.splittingLine_mem_orbit, New.Exchange.thetaTerm_lineRadical}
\lean{New.Exchange.LineIndex, New.Exchange.frameAt_lineTermSplitting, New.Exchange.iCoord_levi_iAdelic, New.Exchange.isAdapted_lineTermSplitting, New.Exchange.lineLevi, New.Exchange.lineLevi_inv_mul_kAdelic_mul_levi, New.Exchange.lineLevi_inv_mul_lAdelic_mul_levi, New.Exchange.lineRadical, New.Exchange.lineTermBasis, New.Exchange.lineTermSplitting, New.Exchange.lineTerm_eq_comp_zeroTerm, New.Exchange.lineTerm_eq_slice, New.Exchange.lineTerm_pairWeil_levi, New.Exchange.planeModel, New.Exchange.planeSplitting, New.Exchange.planeSplitting_iAdelic_mul_planeModel_lAdelic, New.Exchange.tsum_lineTerm}
\lean{New.Exchange.continuous_coneIntegrand, New.Exchange.continuous_ideleAbs_realIdele, New.Exchange.exists_fockMoment_lineTermReal_le, New.Exchange.exists_integrable_fockMoment_coneIntegrand_le, New.Exchange.exists_prod_norm_sq_rpow_le, New.Exchange.exists_sq_rpow_mul_one_add_log_sq_le, New.Exchange.fockMoment_coneIntegrand, New.Exchange.ideleAbs_realIdele_eq_prod, New.Exchange.integrable_coneIntegrand, New.Exchange.isSchwartzVector_integral_coneIntegrand, New.Exchange.realIdele}
\lean{New.Exchange.LineValueFactors, New.Exchange.analyticOnNhd_integral_zetaIntegrand, New.Exchange.baseChange_mul_lineValue_eq_coneKernel_mul, New.Exchange.coneKernel, New.Exchange.coneKernel_eq, New.Exchange.continuousOn_integral_coneKernel, New.Exchange.exists_lineValue_eq, New.Exchange.exists_planeConstantTerm_eq_integral, New.Exchange.exists_planeConstantTerm_eq_thetaLarge, New.Exchange.exists_planeConstantTerm_eq_thetaLarge_cone, New.Exchange.exists_planeConstantTerm_eq_thetaLarge_dualCone, New.Exchange.exists_zetaIntegral, New.Exchange.homSlice_zero_planeDatum, New.Exchange.integrable_coneKernel, New.Exchange.integrable_theta_cone, New.Exchange.integral_coneKernel, New.Exchange.integral_stdKernelForm, New.Exchange.integral_zetaIntegrand_of_one_le, New.Exchange.integral_zetaIntegrand_of_one_lt, New.Exchange.lineValue, New.Exchange.lineValue_zero_of_factor, New.Exchange.pairWeil_stdRowDatum_eq_self, New.Exchange.rankTwist_eq_mul_of_splitting, New.Exchange.shiftIdele, New.Exchange.smul_localDatum, New.Exchange.stdKernelForm, New.Exchange.stdKernelForm_apply, New.Exchange.theta_lineTerm_levi, New.Exchange.tsum_lineValue_poisson, New.Exchange.zetaIntegrand, New.Exchange.zetaIntegrand_eq_tsum, New.Exchange.zetaMeasure}
\lean{Membership.mem.transpose_mem_schwartzBruhat, New.Exchange.IsEquivariantLift.apply_shear, New.Exchange.IsEquivariantLift.centreCoeff_shear, New.Exchange.IsEquivariantLift.comp_mul_mem_ker, New.Exchange.IsEquivariantLift.components_of_mat_eq_one_add, New.Exchange.IsEquivariantLift.exists_linear_centreCoeff, New.Exchange.exists_lineTerm_localDatum, New.Exchange.frameMatrix, New.Exchange.isPrimitiveMatrix_frameMatrix, New.Exchange.ite_isIntegralUnit_eq_ite_invOnUnits_div, New.Exchange.map_shear_stdColDatum_eq, New.Exchange.mul_stdColDatum_triangularForm, New.Exchange.stdColDatum_one_mul_mul, New.Exchange.stdRowDatum_frameAt_lineTermSplitting, New.Exchange.triangularForm, New.areMatched_stdRowDatum_stdColDatum, New.invOnUnits_eq_mul_of_unit, New.isIntMatrix_of_stdColDatum_ne_zero, New.isPrimitiveMatrix_iff_int_and_unit_det, New.isPrimitiveMatrix_of_stdColDatum_ne_zero, New.schwartzBruhat.mem_of_support_isPrimitiveMatrix, New.stdColDatum, New.stdColDatum_mem_schwartzBruhat, New.stdRowDatum, New.stdRowDatum_mem}
\lean{New.Exchange.coneIntegrand_tateVector, New.Exchange.exists_framePoint_eq_of_posProj_eq, New.Exchange.frameAt_real_mul_conjTranspose, New.Exchange.framePoint, New.Exchange.framePoint_apply, New.Exchange.framePoint_eq_mulVec, New.Exchange.framePoint_lineTermSplitting_diag, New.Exchange.framePoint_sub, New.Exchange.isSchwartzVector_tateVector, New.Exchange.lineTermReal_tateVector, New.Exchange.posProj_eq, New.Exchange.posProj_mulVec_framePoint, New.Exchange.realBoxProd_realHeisenbergOp_vacVector, New.Exchange.realPairWeil_tateVector, New.Exchange.signMatrix_mulVec_framePoint, New.Exchange.signMatrix_mul_frameAt_real, New.Exchange.slice_polyVector_apply, New.Exchange.slice_realHeisenbergOp_vacElem_apply, New.Exchange.star_framePoint_dotProduct, New.Exchange.tateVector}
\lean{New.Exchange.analyticAt_ideleAbs_rpow_mul_prod_rpow, New.Exchange.analyticOnNhd_gammaIntegral, New.Exchange.exists_analyticOnNhd_eq_L_mul_gammaIntegral, New.Exchange.exists_analytic_extension_functional_equation, New.Exchange.exists_analytic_extension_ne_zero, New.Exchange.exists_exp_neg_sq_mul_rpow_mul_one_add_log_sq_le, New.Exchange.gammaIntegral, New.Exchange.gammaIntegral_pos, New.Exchange.gammaWeight, New.Exchange.gammaWeight_pos, New.Exchange.integrable_gammaWeight_mul_rpow, New.Exchange.integrable_prod_exp_neg_sq_mul_rpow, New.Exchange.integral_coneIntegrand_tateVector, New.Exchange.pairTheta_source_zero, New.Exchange.pairTheta_target_zero, New.Exchange.schrodingerOp_vacElem_of_isEmpty, New.Exchange.stdColDatum_one_one, New.Exchange.thetaLarge_smul, New.Exchange.thetaLarge_stdSpace_empty_vacVector_one, New.Exchange.theta_vacElem_of_isEmpty}
\lean{New.integrable_prod_inv_one_add_log_sq, New.piUnits_val_apply_ne_zero, New.principalAdeleUnits}

\begin{proof}[Sources]
For cuspidal forms the constant term of a theta lift along an isotropic line is the lift for the pair with one plane removed; this is Rallis' tower property \cite[Thm.~I.1.1]{RallisHowe1984},
and for unitary groups Wu \cite[Prop.~3.1]{Wu2013}. Here the lifted function is a character, the data are ramified and the isotropic orbit survives, as a Tate integral \cite{Tate1967}; the Lean development
follows the classical route. (a) The Euler product converges absolutely for $x>1$ and has no zero at $x=1$ because a character of nonzero weight is not trivial; the recursion is the definition of
$L_*$. (b) The line term decays rapidly in $a$ in every Fock moment, while $\abs a^z$ grows at most polynomially, so the integrand is dominated by a continuous function of $\kappa$ on a compact group
times $\prod_w(1+\log^2\abs{a_w})^{-1}$, which is integrable for Haar measure on $(\C^\times)^{\#w}$. (c) The constant term of the theta kernel is the sum over the rational vectors of $V'$ of their terms;
a term vanishes unless the vector is isotropic, and the term of $0$ vanishes on the plane datum, since the standard datum lives on matrices of full rank. By Witt's theorem the nonzero isotropic lines
form one orbit of the rational points, so the integral over $[\U(V')]$ unfolds to the stabiliser of the line of $X$ and then, in Iwasawa coordinates $((u',h),\kappa)$ with the module
$\abs{u'}^{\operatorname{rk}V+1}$, to an integral over the idele classes, $[\U(V)]$ and the maximal compact subgroup. Levi points act on the line terms by a scalar and a rescaling of the adele, which turns
the integrand into Tate's integrand $\nu(a/\bar a)\abs a^{x}\sum_{\xi\in E^\times}$ of line values; Poisson's formula on the line identifies the two models. At $S$ the datum factors along the line into
a Gauss sum of modulus $(\prod_Sq_w)^{-1/2}$ times the finite factor $a\mapsto\chi^{-1}(a)\chi(\varpi\bar a)$, and Tate's finite integral of this factor is $L^S(x,\nu)$ times a volume for $x>1$. Neither side
of Poisson's formula has a term at $0$, so the integral is a Mellin transform convergent on every strip and analytic on the real line; continuity in the exponent carries the formula to $x=1$. (d) is
(c) at $x=z+1$, on the side of Poisson's formula fixed by the sign of $e$; the modulus of the Gauss sum and of $\abs{\alpha_0}^{z+1}$ go into $C$. (e) Apply (c) to the pair $(H,H)$ over the space $0$ with
the test vector, the Fock vector of the product of the $\mu(w)$-th powers of the diagonal variables: the maximal compact subgroup fixes it, and its cone integrand is the vacuum times
$e^{-\sum_w\abs{a_w}^2}\abs a^x\prod_w\abs{a_w}^{2\mu(w)}$, whose integral is a positive analytic function of $x>-k$. Dividing $\mathcal Z$ by the analytic factor without zeros that multiplies $L^S$
continues $L^S$ to $x>-k$, and comparing the two sides of Poisson's formula gives the functional equation in absolute values \cite[\S4.4]{Tate1967}; the continuation has no zero at an integer $j>-k$
because, by the functional equation, its value there is a nonzero multiple of $L^\sharp(1-j)$ with $1-j\ge1$ when $j\le0$, and is an Euler value at $j\ge1$.
\end{proof}

\subsection{Vanishing descents and periods}\label{ss:B-periods}
This subsection supplies the passage from vanishing theta descents to vanishing periods used for the false lifts (\cref{thm:false-lift-periods}). Let $W$, $V$ be hermitian spaces, $a$ a character of $\U(1)(\Ad)$,
$S$ a set of finite places and $Y$ an automorphic function on $\U(V)(\Ad)$. The descents of $Y$ to $W$ \emph{vanish} (for $a$ and $S$) if $g\mapsto\int_{[\U(V)]}\theta(\omega(g,h)f)\,a(\det h)\,Y(h)\,dh$ is zero for every datum
$f$ on $\Hom(W,V)$ that is the vacuum at the places of $S$ and their conjugates. If $W$ has a splitting with isotropic frames of rank $p$, the theta kernel at $(g,h)$ is the sum of the \emph{terms} of the rational frames
$x$ (the $p$-tuples of vectors of $V$), and the \emph{Fourier coefficient} at an hermitian $p\times p$ matrix $\beta$ of a function $F$ on $\U(W)(\Ad)$ is the mean of $\psi(\operatorname{tr}(j\beta M))F(z(M)g)$ over the adelic $M$
modulo the rational ones, $z$ parametrising the centre of the unipotent radical, $\psi$ the standard additive character and $j$ the purely imaginary element of the Weil data. For a splitting $V=V''\oplus V^\circ$ of
the definite space of rank $n$ into blocks of ranks $n_1$, $n_2$ and a point $h_1$ of $\U(V'')(\Ad)$, the \emph{block period} of $Y$ is $I(h)=\int_{[\U(V^\circ)]}Y(\mathrm{diag}(h_1,h_2)h)\,dh_2$.

\begin{theorem}[Vanishing descents and periods]\label{thm:B-periods}
\begin{enumerate}[label=\textup{(\alph*)}]
\item (Testing and translation.) Let $V$ be definite and $A$ a set of finite places disjoint from its conjugates. The descents of $Y$ vanish for the data that are the vacuum off $A$ and its conjugates as soon as
they vanish for the data $v_P\otimes\Psi\otimes\mathrm{vac}$, $P$ a polynomial and $\Psi$ a Schwartz--Bruhat function above $A$. If the descents of $Y$ vanish for $S$, so do those of its right translate by a point that
is integral at every place of $S$.
\item (Fourier coefficients.) Let $y$ be continuous on $\U(V)(\Ad)$ and left invariant under the rational points, with terms of the kernel at $g$ absolutely integrable against $y$. The Fourier coefficient at $j\beta$
of the theta lift of $y$ is the sum, over the rational frames with Gram matrix $\beta$, of the integrals over $[\U(V)]$ of their terms times $y$. If $\beta$ is invertible and $x_0$ is one such frame, it is the integral
of the term of $x_0$ times $y$ modulo the rational stabiliser of $x_0$.
\item (Periods.) Let $T$ be a set of finite places and $Y$ continuous on $[\U(V)]$, $V=V''\oplus V^\circ$ definite, right invariant under the integral points with component $1$ off $T$ and its conjugates. If the descents
of $Y$ to the split hermitian space of rank $2n_1$, twisted by $\upsilon^{n_1}$, vanish for all data that are the vacuum at the places of $T$ and their conjugates, then the block period $I(1)$ vanishes for every $h_1$.
\end{enumerate}
\end{theorem}
\lean{New.SiegelWeil.DescentsVanish, New.SiegelWeil.DescentsVanish.comp_mul_right, New.SiegelWeil.DescentsVanish.integral_frameChar_eq_zero, New.SiegelWeil.DescentsVanish.integral_frameChar_mul_blockPeriod_eq_zero, New.SiegelWeil.DescentsVanish.of_schrodingerOp, New.SiegelWeil.blockEmb_eq_secondBlock_mul, New.SiegelWeil.blockEmb_mul_firstFrame, New.SiegelWeil.blockPeriod, New.SiegelWeil.blockPeriod_eq_of_frameChar_eq, New.SiegelWeil.blockPeriod_eq_of_movedFrame_eq, New.SiegelWeil.blockQuotMeasure, New.SiegelWeil.comap_Gamma_secondBlock, New.SiegelWeil.continuous_blockPeriod, New.SiegelWeil.continuous_frameChar, New.SiegelWeil.continuous_frameWeight, New.SiegelWeil.continuous_movedFrame, New.SiegelWeil.exists_continuousMap_comp_eq, New.SiegelWeil.exists_continuousMap_eq_conj_blockPeriod, New.SiegelWeil.exists_movedFrame_eq, New.SiegelWeil.exists_theta_homSlice_eq_mul_frameChar, New.SiegelWeil.firstFrame, New.SiegelWeil.frameChar, New.SiegelWeil.frameWeight, New.SiegelWeil.frameWeight_ne_zero_iff, New.SiegelWeil.hasIntegrableTerms_hyperbolicSplitting_of_isRapid, New.SiegelWeil.integral_blockEmb_eq_zero_of_descentsVanish, New.SiegelWeil.integral_continuousMap_frameChar_eq_zero, New.SiegelWeil.inv_blockEmb_mul_secondBlock_mul, New.SiegelWeil.isClosedEmbedding_secondBlock, New.SiegelWeil.mem_map_Gamma_secondBlock, New.SiegelWeil.mem_range_secondBlock, New.SiegelWeil.movedFrame, New.SiegelWeil.movedFrame_inv_blockEmb_mul, New.SiegelWeil.movedFrame_secondBlock_mul, New.SiegelWeil.secondBlock, New.SiegelWeil.smul_eq_of_frameChar_eq, New.Space.Splitting.centreCoeff_eq_integral_ratStabilizer, New.Space.Splitting.centreCoeff_eq_integral_stabilizer, New.Space.Splitting.centreCoeff_integral_eq_tsum, New.Space.Splitting.fixesLeft_thetaTerm_mul, New.Space.Splitting.hasSum_thetaTerm, New.Space.Splitting.integral_pairTheta_radicalCentre_mul, New.Space.Splitting.radicalCentre_apply, New.Space.Splitting.thetaTerm_mul_ratEmbed_mul, New.Space.Splitting.thetaTerm_radicalCentre_mul, New.Space.gram_eq_iff_mem_orbit, New.Space.mem_ratStabilizer, New.Space.ratStabilizer, New.U.smul_matrix_eq_mul}

\begin{proof}[Sources]
Fourier coefficients of theta lifts are periods over stabilisers, and vanishing lifts in a theta tower kill periods on the complementary group; see Rallis \cite{RallisHowe1984} and, for this unfolding,
Walls \cite[Prop.~2.1.1]{Walls2015}. Here a single nondegenerate coefficient of a lift to a split unitary group is unfolded once. (a) The descent at $g$ is a bounded linear functional of the datum, so it suffices to
test the data $v_P\otimes\varrho(x)\mathrm{vac}$ with $x$ finite and supported on $A$ and its conjugates, and $\varrho(x)\mathrm{vac}$ is a multiple of the image of the vacuum under the Heisenberg operators realising a
Schwartz--Bruhat function of the components above $A$. The descent of the translate of $Y$ by $h_0$ by a datum $f$ is the descent of $Y$ by $\omega(1,h_0^{-1})f$, again the vacuum at $S$ as $h_0$ is integral there. (b)
Left multiplication by the central element $z(M)$ multiplies the term of $x$ by $\psi(\operatorname{tr}(-j\beta(x)M))$, $\beta(x)$ the Gram matrix of $x$, and for $y$ invariant under the rational points the term of $x$ times $y$ is
unchanged under $(x,h)\mapsto(\gamma x,\gamma h)$. Expanding the kernel into its terms and exchanging the sum with the integral and the mean, the mean over $M$ keeps the frames with $\beta(x)=\beta$. For invertible $\beta$,
Witt's theorem makes these frames one orbit of the rational unitary group, and the sum over the orbit unfolds to the stabiliser. (c) Let $o_1$ be the frame of the first block, $\alpha_0$ the scaling idele of the Weil data,
$\mathrm{vac}[x]=\abs{e(\alpha_0x)}^2$ the slice of the vacuum and $\psi_u(x)=\psi(\delta\sum u_{ij}x_{ij})$, with $\delta$ generating the inverse different and the components at $T$ and its conjugates removed. For the Heisenberg
translate of the vacuum along the isotropic vector with coordinates $\alpha_0\bar u$, which is the vacuum at $T$, the term of the frame $o_1$ at the Levi point of $\alpha_0h_1$ is $c\,\psi_u(x)\mathrm{vac}[x]$ with
$x=g^{-1}o_1h_1$ and $c\ne0$; the theta integral against the continuous $Y$ on the compact quotient converges term by term. By (b) with $\beta=1$ the vanishing descent gives $\int\psi_u(x)\mathrm{vac}[x]Y(g)\,dg=0$ modulo the
rational points of the second block, which is the stabiliser of $o_1$; integrating first over $[\U(V^\circ)]$ gives $\int\psi_u(h^{-1}o_1)\mathrm{vac}[h^{-1}o_1]I(h)\,dh=0$ over $\U(V^\circ)(\Ad)\backslash\U(V)(\Ad)$, for every $u$. The functions
that agree with some $\psi_u$ on frames form a monoid closed under conjugation that separates frames up to their components at $T$, so by the Stone--Weierstrass theorem the same holds with $\psi_u$ replaced by any
continuous function $\chi$ of the values $u\mapsto\psi_u(h^{-1}o_1)$. Where $\mathrm{vac}[h^{-1}o_1]\ne0$ the frame is integral, so it is $\gamma^{-1}o_1$ for an integral $\gamma$ by Witt's theorem over $\hat\OO$, and $I(h)$ depends only on the
values $\psi_u(h^{-1}o_1)$, because the components at $T$ are absorbed by the integral points supported there, which fix $Y$; by Tietze's theorem $\bar I$ is such a $\chi$. Hence $\int\mathrm{vac}[h^{-1}o_1]\abs{I(h)}^2\,dh=0$, so $I$
vanishes where $\mathrm{vac}\ne0$, in particular at $h=1$, where the frame is $o_1$.
\end{proof}

\subsection{Rallis' inner product formula}\label{ss:B-rallis}
Let $S$, $T$ be hermitian spaces of ranks $n$, $m$ with $2n<m$, and let $\omega$ be the twisted Weil representation of $\U(S)\times\U(T)$ on the data $\mathcal S(\Hom(S,T))$. The \emph{doubled
datum} of $F$ lives on $\Hom(S\oplus S^-,T)$. A finite datum $F$ is \emph{clean} if only the Lagrangians of $S\oplus S^-$ that are graphs of isometries contribute to its Siegel--Weil
formula, and \emph{sharp} for an open subgroup $J$ off a set $S_0$ of places if it is the vacuum off $S_0$ and, at the finite points $g$ integral off $S_0$, $\langle F,\omega(g)F\rangle$ is
$\langle F,F\rangle$ on $J\cap K_0$ and $0$ otherwise; its \emph{mass} $c$ is $\langle F,F\rangle/[K_0:J\cap K_0]$. The \emph{vacuum coefficient} of a finite point $\gamma$ supported off $S_0$ is
$\mathrm{mc}(\gamma)=(b(\det\gamma)\,\mathrm{fm}(\gamma)^m)^{-1}$, the factor by which $\gamma$ changes the matrix coefficient of the vacuum, and the \emph{Rallis series} is
$\sum_\gamma\mathrm{mc}(\gamma)R(\gamma)$ over the cosets of the points integral off $S_0$.
\lean{New.SiegelWeil.negSourceFactor, New.SiegelWeil.norm_negSourceFactor, New.SiegelWeil.theta_negSourceFactor_smul, New.SiegelWeil.doubleDatum, New.SiegelWeil.graphSplitting, New.SiegelWeil.cosetTerm, New.SiegelWeil.weilMatrixCoeff, New.SiegelWeil.IsClean, New.SiegelWeil.IsSharp, New.SiegelWeil.doubleDatum_apply, New.SiegelWeil.theta_doubleDatum, New.SiegelWeil.rho_fromCols_boxProd, New.SiegelWeil.doubleDatum_apply_fromCols, New.SiegelWeil.rho_fromCols_zero_doubleDatum, New.SiegelWeil.zeroTerm_graphSplitting, New.SiegelWeil.zeroTerm_doubleDatum_rho, New.SiegelWeil.pairWeil_rho_shiftMatrix, New.SiegelWeil.mul_eq_zero_of_zeroTerm_ne_zero, New.SiegelWeil.exists_graph_eq_of_zeroTerm_ne_zero, New.SiegelWeil.graphSplitting_lagrangian_injective, New.SiegelWeil.theta_splitWeil_doubleDatum, New.SiegelWeil.zeroTerm_graphSplitting_doubleDatum, New.SiegelWeil.integral_conj_theta_mul_theta, New.SiegelWeil.exists_doubling_bound, New.SiegelWeil.exists_doubling_le_height, New.SiegelWeil.measurePreserving_quotMeasure_unfold, New.SiegelWeil.enorm_mul_height_pow_le, New.SiegelWeil.inner_product_formula, New.SiegelWeil.weilMatrixCoeff_mul, New.SiegelWeil.exists_integral_isotropic_frame}

\begin{theorem}[Rallis' inner product formula]\label{thm:B-rallis}
Let $2n<m$, and let the real vector $v$ be an edge from a type $e$ (for instance the joint harmonic of corresponding parameters).
\begin{enumerate}[label=\textup{(\alph*)}]
\item (Doubling.) For clean data, the integral
\[
  \int_{[\U(T)]}\overline{\theta(\omega(g,h)\phi)}\,\theta(\omega(g',h)\phi')\,dh
\]
is the sum of the zero terms of the doubled datum along the graphs of isometries.
\item Let $F$ be sharp for $J$ off $S_0$ with mass $c$, let $z$ be the Rallis series of $y$ and $\Omega$ the real Weil representation. Then $\langle\Theta y',\Theta y\rangle$ is $\vol[\U(T)]\cdot c$ times the absolutely
convergent integral $\int\langle v,\Omega(x)v\rangle\langle y',R(x)z\rangle\,dx$ over the real points $x$ of $\U(S)$.
\item There is $A\ne0$, depending only on the spaces and the real data, such that for $F$ sharp for $J$ off $S_0$ with mass $c$, the lift by $v\otimes F$ is injective on the $J$-invariant forms of
type $e$, and $\langle\Theta y',\Theta y\rangle=A\,c\sum_\gamma\mathrm{mc}(\gamma)\langle y',R(\gamma)y\rangle$.
\item If the signs of $S$ agree or $\mathfrak m\ne(1)$, a function $\Psi\ne0$ on matrices modulo $\mathfrak m$ that is fixed with the vacuum by a level $J(\mathfrak m)\le J\le K_0$, and whose support is
carried into $J$ by a fixed matrix, gives a datum $\Psi\otimes\mathrm{vac}$ sharp for $J$ off the places of $\mathfrak m\bar{\mathfrak m}$, with mass the mean of $\abs\Psi^2$ over $[K_0:J]$.
\item (Rank one.) For $n=1$, $m=2k$, $k>1$, and a character $\nu$ with tame part trivial off $S$, the sum of the vacuum coefficients of $\nu\upsilon^k$ over the cosets off $D=S\cup\bar S$ is
$L^D(k,\nu^{-1})$ divided by the Euler product at $k$ of the place signs off $D$, a partial value of $\zeta_F(2k)$.
\end{enumerate}
\end{theorem}
\lean{New.SiegelWeil.IsSharp.isOpen_level, New.SiegelWeil.IsSharp.mass_ne_zero, New.SiegelWeil.IsSharp.integral_ite_eq_mass, New.SiegelWeil.weilMatrixCoeff_integrated, New.SiegelWeil.IsSharp.weilMatrixCoeff_integrated_mul, New.SiegelWeil.IsSharp.inf_supportedOn_le, New.SiegelWeil.IsSharp.isUnramifiedOn, New.hasSum_setIntegral_mul_coset, New.SiegelWeil.IsSharp.hasSum_integral_weilMatrixCoeff_mul, New.SiegelWeil.IsSharp.petersson_lift_eq_integral, New.SiegelWeil.realConv_eq_zero_of_lift_eq_zero, New.SiegelWeil.exists_petersson_formLift_eq_rallisSeries, New.SiegelWeil.rallis_inner_product_cosetTerm, New.SiegelWeil.IsSharpFn, New.SiegelWeil.IsSharpFn.isPrimitiveMatrix, New.SiegelWeil.IsSharpFn.eq_zero_or_eq_zero, New.SiegelWeil.pairWeil_stdLevel_vacuumDatum, New.SiegelWeil.IsSharpFn.weilMatrixCoeff_eq_zero_of_mem_stdLevel_of_notMem, New.SiegelWeil.IsSharpFn.weilMatrixCoeff_eq_zero_of_notMem_stdLevel, New.SiegelWeil.IsSharpFn.isSharp, New.SiegelWeil.isSharpFn_of_mul_eq_one, New.SiegelWeil.isSharp_vacuumDatum_of_mul_eq_one, New.AutChar.HasWeight.apply_eq_one, New.SiegelWeil.vacuumCoeff_eq_spaceChar, New.placeSign, New.norm_placeSign_le, New.placeSign_factor_mul, New.glMultiplier_singleGL_uniformizer_zpow, New.mem_localIntGL_unit_iff, New.vacuumCoeff_ofSplitPlace_uniformizer_zpow, New.exists_hasSum_vacuumCoeff_at_place, New.vacuumSum_pair, New.vacuumSum_eq_L_div_eulerProd}

\begin{proof}[Sources]
This is Rallis' inner product formula by doubling \cite{Rallis1987}, in Weil's range and for rapidly decreasing forms; the unitary unfolding is as in Li \cite{Li1992}. (a) applies the
Siegel--Weil formula (\cref{thm:B-siegel-weil}) to $S\oplus S^-$; for a clean datum the boundary orbits drop out. (b) integrates (a) against two functions on $[\U(S)]$ and, for a sharp datum,
cuts the finite points into cosets, on each of which the matrix coefficient is the vacuum coefficient. (c) The injectivity is Rallis' argument \cite{Rallis1984}: if the lift vanishes, the real
matrix coefficient of $v$ convolves the form to zero, and for an edge this convolution is a nonzero multiple of the form. (d) The matrix coefficient of $\Psi\otimes\mathrm{vac}$ at an integral point is
a sum of products of values of $\Psi$, which vanishes off $J$ by the support condition. (e) At a split place the vacuum coefficients form a geometric series in $\nu(\varpi)q^{-k}$, and at the other
places the coset is trivial; the product over the places is the partial $L$-value over the partial Euler product.
\end{proof}
